\section{The W(3,3) Uniqueness Theorem}
\label{sec:uniqueness}

Let $q$ be a prime power and $W(3,q)$ the Lubotzky--Phillips--Sarnak
Ramanujan bipartite graph with degree $k=q(q+1)$,
vertex count $v=(q+1)(q^2+1)$ per side, and adjacency eigenvalues
$r = q-1$ and $s = -(q+1)$, with multiplicities
$f = q(q+1)^2/2$ and $g = q(q^2+1)/2$.
Define
\[
  \Phi_3(q) = q^2+q+1,\qquad
  \Phi_4(q) = q^2+1,\qquad
  \Phi_6(q) = q^2-q+1.
\]
The Ihara zeta function factors through the two quadratic polynomials
\[
  p_1(u) = 1 - (q-1)u + (k-1)u^2,\qquad
  p_2(u) = 1 + (q+1)u + (k-1)u^2,
\]
whose shared constant term is $k(q)-1 = q^2+q-1$.

\begin{definition}[Spectral deficit]
Define the spectral deficit by
\[
  \mathrm{Deficit}(q) = -\frac{\mathrm{disc}(p_1)}{4} - \Phi_4(q).
\]
Since
\[
  \mathrm{disc}(p_1) = (q-1)^2 - 4(k-1) = -3q^2-6q+5,
\]
we obtain the exact factorization below.
\end{definition}

\begin{lemma}[C1]
\[
  \mathrm{Deficit}(q) = -\frac{(q-3)^2}{4}.
\]
Hence $\mathrm{Deficit}(q)=0$ if and only if $q=3$.
\end{lemma}
\begin{proof}
A direct expansion gives
\[
  \mathrm{Deficit}(q)
  = \frac{3q^2+6q-5}{4} - (q^2+1)
  = \frac{-q^2+6q-9}{4}
  = -\frac{(q-3)^2}{4}. \qedhere
\]
\end{proof}

\begin{lemma}[C2]
$k(q)+g(q)=q^q$ if and only if $q=3$.
\end{lemma}
\begin{proof}
We compute
\[
  k(q)+g(q)=q(q+1)+\frac{q(q^2+1)}{2} = \frac{q^3}{2}+q^2+\frac{3q}{2}.
\]
At $q=3$ this equals $27=3^3$.
For $q\ge 4$, the right-hand side $q^q$ dominates any cubic polynomial,
so no further positive integer solutions exist.
\end{proof}

\begin{lemma}[C3]
The identities
\[
  -f(q)=\tau(2),\qquad k(q)\,q\,\Phi_6(q)=\tau(3)
\]
hold if and only if $q=3$.
\end{lemma}
\begin{proof}
Since $\tau(2)=-24$ and $\tau(3)=252$, we evaluate at $q=3$:
\[
  -f(3)=-\frac{3\cdot 4^2}{2}=-24,
  \qquad
  k(3)\cdot 3\cdot \Phi_6(3)=12\cdot 3\cdot 7=252.
\]
The first equation has unique positive integer solution $q=3$,
and the second is strictly increasing for $q\ge 2$.
\end{proof}

\begin{lemma}[C4]
$\mathbb{Q}(\sqrt{-\Phi_6(q)})$ has class number $1$ for $q\in\{2,3,7\}$.
In conjunction with C1, this is unique to $q=3$.
\end{lemma}
\begin{proof}
The Heegner discriminants are
\(
\{1,2,3,4,7,8,11,19,43,67,163\}.
\)
Now
\[
  \Phi_6(2)=3,\qquad \Phi_6(3)=7,\qquad \Phi_6(7)=43,
\]
all of which are Heegner. But C1 fails at $q=2$ and $q=7$.
\end{proof}

\begin{lemma}[C5]
Let $E_{-\Phi_6(q)}$ be the CM elliptic curve attached to the order of
 discriminant $-\Phi_6(q)$. Then
\[
  a_{k(q)-1}(E_{-\Phi_6(q)}) = \mathrm{ev}_s(q)=-(q+1)
\]
holds at $q=3$ and fails already at $q=2$.
\end{lemma}
\begin{proof}
At $q=3$, one has $k-1=11$, $\Phi_6=7$, and
\[
  a_{11}(E_{-7})=-4=-(3+1)=\mathrm{ev}_s(3).
\]
At $q=2$, one has $k-1=5$, $\Phi_6=3$, while direct point-counting gives
\[
  a_5(E_{-3})=0\neq -3=\mathrm{ev}_s(2).
\]
So the identity is not generic in the $W(3,q)$ family.
\end{proof}

\begin{theorem}[W(3,3) Uniqueness]
$W(3,3)$ is the unique member of the $W(3,q)$ family satisfying
all five characterizations C1--C5 simultaneously.
\end{theorem}
\begin{proof}
By the preceding lemmas, C1, C2, and C3 each isolate $q=3$,
while C1$\wedge$C4 and C5 also hold uniquely at $q=3$.
Therefore the full conjunction C1$\wedge$C2$\wedge$C3$\wedge$C4$\wedge$C5
holds if and only if $q=3$.
\end{proof}
